\documentclass[11pt]{article}
\usepackage[margin=1in]{geometry}
\usepackage[T1]{fontenc}
\usepackage[utf8]{inputenc}
\usepackage{amsmath,amssymb,amsthm}
\usepackage{enumitem}

\title{ICPC World Finals 2023\\K. Alea Iacta Est}
\author{}
\date{}

\begin{document}
\maketitle

\section*{Problem Summary}

We repeatedly roll $d$ dice. After each roll we may keep any subset of the dice and reroll the rest. A
roll is successful if the top letters can be permuted to form some dictionary word. We must compute the
minimum possible expected number of rolls until success, or decide that success is impossible.

\section*{Key Observations}

\begin{itemize}[leftmargin=*]
    \item A state is a vector in $\{1,2,3,4,5,6,*\}^d$. A fixed state has no \texttt{*} symbols and describes
    one concrete roll outcome. An undetermined state has stars in the positions we are about to reroll.
    \item The goal states are exactly the fixed states whose letters, after sorting, match some sorted
    dictionary word.
    \item From a fixed state we may, at cost $0$, choose any nonempty subset of coordinates and replace
    them by stars. That moves us to an undetermined state.
    \item If an undetermined state has $s$ stars, then one reroll replaces them uniformly by one of $6^s$
    matching fixed states and costs $1$ roll.
    \item This yields Bellman equations on a finite state graph. Running the relaxations \emph{backwards}
    from goal states leads to a Dijkstra-like algorithm.
\end{itemize}

\section*{Algorithm}

\begin{enumerate}[leftmargin=*]
    \item Encode every state in base $7$, using digit $6$ for \texttt{*}. Precompute for each state whether it
    is fixed and how many stars it contains.
    \item Initialize distance $0$ for every goal fixed state and $\infty$ for all other states. Process states in
    increasing tentative distance with a priority queue.
    \item When a fixed state $x$ is finalized, update every undetermined state $y$ obtained by replacing a
    nonempty subset of coordinates of $x$ by stars.
    \item For one such $y$, let $D_y$ be the set of matching fixed states that have already been finalized. If
    $y$ has $s$ stars, then using only those finalized fixed states as stopping states gives candidate value
    \[
    \frac{6^s + \sum_{z \in D_y} \mathrm{dist}[z]}{|D_y|}.
    \]
    Maintain this value incrementally by storing $|D_y|$ and the sum of finalized distances.
    \item When an undetermined state $y$ is finalized, update every matching fixed state $x$ with
    \[
    \mathrm{dist}[x] \le \mathrm{dist}[y],
    \]
    because from $x$ we may immediately choose to reroll exactly the starred coordinates of $y$.
    \item The answer is the distance of the all-star state. If it stays infinite, output \texttt{impossible}.
\end{enumerate}

\section*{Correctness Proof}

We prove that the algorithm returns the correct answer.

\paragraph{Lemma 1.}
For a fixed state $x$ and any nonempty subset of coordinates, moving from $x$ to the corresponding
undetermined state $y$ has additional expected cost $0$.

\paragraph{Proof.}
Choosing which dice to reroll is a free decision taken after observing the current roll. No new roll is
performed yet, so the expected number of \emph{future} rolls from $x$ after choosing that subset is exactly
the expected number of future rolls from $y$. \qed

\paragraph{Lemma 2.}
Let $y$ be an undetermined state with $s$ stars, and let $D$ be any nonempty set of matching fixed states.
If we keep rerolling those starred coordinates until the outcome first lands in $D$, then the expected
remaining number of rolls is
\[
\frac{6^s + \sum_{z \in D}\mathrm{dist}[z]}{|D|}.
\]

\paragraph{Proof.}
Each roll of the $s$ starred dice is uniformly distributed over the $6^s$ matching fixed states, so the
probability of landing in $D$ on one attempt is $|D|/6^s$. Therefore the expected number of attempts
until the first success is $6^s/|D|$.

Conditioned on success, the reached state is uniformly distributed over $D$, so the expected future cost
after that success is
\[
\frac{1}{|D|}\sum_{z \in D}\mathrm{dist}[z].
\]
Adding these two expectations gives the stated formula. \qed

\paragraph{Lemma 3.}
When the algorithm finalizes a state, its recorded distance equals the true optimal expected number of
future rolls from that state.

\paragraph{Proof.}
We process states in nondecreasing tentative distance. For fixed states, every relaxation to them comes
from an already finalized undetermined state by Lemma 1, so their tentative value is the best value
obtainable from already optimal successor states.

For an undetermined state $y$, Lemma 2 shows that any strategy based on a set $D$ of already finalized
matching fixed states yields exactly the candidate value maintained by the algorithm. Adding a not-yet
finalized fixed state cannot improve the optimum before that state's own distance is known, because all
finalized states have distance no larger than any unfinalized one. Thus the first time $y$ is extracted from
the priority queue, its tentative value already matches the true optimum. The usual Dijkstra induction now
applies to both kinds of states. \qed

\paragraph{Theorem.}
The algorithm outputs the minimum expected number of rolls needed to form a dictionary word, or
\texttt{impossible} if no word can ever be formed.

\paragraph{Proof.}
By Lemma 3, every finalized distance is optimal. The all-star state represents the situation before the first
roll, so its optimal expected future number of rolls is exactly the answer required by the statement. If its
distance remains infinite, then no goal state is reachable under any strategy, so the correct output is
\texttt{impossible}. \qed

\section*{Complexity Analysis}

There are $7^d$ total states and at most $6^d$ fixed states. Since $d \le 6$, this is small enough to
enumerate explicitly. Each state's relaxations are over subsets or completions of at most $d$ coordinates,
so the total running time is polynomial in $7^d$ and easily fits the limit. The memory usage is $O(7^d)$.

\section*{Implementation Notes}

\begin{itemize}[leftmargin=*]
    \item Dictionary words are stored in sorted form, because the order of dice does not matter when
    checking whether a fixed state forms a word.
    \item Distances are stored in floating point and printed with nine digits after the decimal point.
\end{itemize}

\end{document}
